\section{Auxiliary Theorems}
\label{appendix:auxiliary-theorems}

This appendix collects the closed-form propositions used by the Tipping Point method of Section~\ref{sec:tipping-point}.

We begin with an elementary identity used repeatedly below.

\begin{lemma}[Weighted-mean identity]\label{lem:weighted-variance}
Let \(w_1,\dots,w_s\geq 0\) with \(W\defsymbol\sum_{i=1}^s w_i>0\), and let \(x_1,\dots,x_s\in\mathbb{R}\). Set \(\mu\defsymbol\frac{1}{W}\sum_{i=1}^s w_ix_i\). Then
\[
\min_{y\in\mathbb{R}}\sum_{i=1}^s w_i(x_i-y)^2
=\sum_{i=1}^s w_i(x_i-\mu)^2
=\frac{1}{W}\sum_{1\leq i<j\leq s}w_iw_j(x_i-x_j)^2.
\]
\end{lemma}
\begin{proof}
The function \(y\mapsto\sum_i w_i(x_i-y)^2\) is convex with derivative \(-2\sum_i w_i(x_i-y)\), which vanishes at \(y=\mu\); this gives the first equality. For the second,
\[
\sum_i w_i(x_i-\mu)^2
=\sum_i w_ix_i^2-W\mu^2
=\frac{W\sum_i w_ix_i^2-\bigl(\sum_i w_ix_i\bigr)^2}{W},
\]
and symmetrizing in \((i,j)\) the numerator equals
\(\sum_{i,j}w_iw_j(x_i^2-x_ix_j)=\tfrac12\sum_{i,j}w_iw_j(x_i-x_j)^2=\sum_{i<j}w_iw_j(x_i-x_j)^2\).
\end{proof}

\begin{proposition}[General mass transfer to a more extreme point]
\label{prop:mec-general-mass-shift-formulas}
Let \(0\le l<r<r'\le 1\), let \(m,n>0\), and set \(d\defsymbol r-l\), \(e\defsymbol r'-r\), and \(C\defsymbol m^2+n^2\). Let
\[
\pi=((m,n),(l,r)),
\]
and, for \(k\in(0,n]\), let
\[
\pi_{r:k\rightarrow r'}=((m,n-k,k),(l,r,r')).
\]
Then
\[
\MEC[2,2](\pi)=\frac{m^2n^2d^2}{C},
\]
and
\[
\MEC[2,2](\pi_{r:k\rightarrow r'})
=
\frac{
m^2(n-k)^2d^2+m^2k^2(d+e)^2+(n-k)^2k^2e^2
}{
2k^2-2nk+C
}.
\]
Hence, if
\[
\Delta(k)\defsymbol \MEC[2,2](\pi_{r:k\rightarrow r'})-\MEC[2,2](\pi),
\]
then
\[
\Delta(k)=\frac{k\,R_{m,n,d,e}(k)}{Q_{m,n}(k)},
\]
where
\[
R_{m,n,d,e}(k)
=
e^2Ck^3-2e^2nCk^2+\bigl(2d^2m^4+2de\,m^2C+e^2C^2\bigr)k-2d^2m^4n,
\]
and
\[
Q_{m,n}(k)=C(2k^2-2nk+C).
\]
In particular, \(Q_{m,n}(k)>0\) for all \(k\in[0,n]\), so the sign of \(\Delta(k)\) for \(k>0\) is determined by \(R_{m,n,d,e}(k)\).
\end{proposition}

\begin{proof}
By Lem.~\ref{lem:weighted-variance} applied to \(\pi\) with weights \(m^2\) and \(n^2\) at positions \(l\) and \(r\), and using \(r-l=d\),
\[
\MEC[2,2](\pi)=\frac{m^2n^2(r-l)^2}{m^2+n^2}=\frac{m^2n^2d^2}{C}.
\]
Applying the same lemma to \(\pi_{r:k\rightarrow r'}\), whose weights are \(m^2\), \((n-k)^2\), \(k^2\) at positions \(l\), \(r\), \(r'\), with the squared distances \((r-l)^2=d^2\), \((r'-l)^2=(d+e)^2\), \((r'-r)^2=e^2\),
\[
\MEC[2,2](\pi_{r:k\rightarrow r'})
=
\frac{m^2(n-k)^2d^2+m^2k^2(d+e)^2+(n-k)^2k^2e^2}{m^2+(n-k)^2+k^2}.
\]
Since \(m^2+(n-k)^2+k^2=2k^2-2nk+C\), this gives the stated formula for \(\MEC[2,2](\pi_{r:k\rightarrow r'})\).

Set \(Q_{m,n}(k)\defsymbol C(2k^2-2nk+C)\). Putting \(\Delta(k)=\MEC[2,2](\pi_{r:k\rightarrow r'})-\MEC[2,2](\pi)\) over this common denominator,
\[
\Delta(k)
=
\frac{C\bigl[m^2(n-k)^2d^2+m^2k^2(d+e)^2+(n-k)^2k^2e^2\bigr]-m^2n^2d^2(2k^2-2nk+C)}{Q_{m,n}(k)}.
\]
At \(k=0\) the second weight reduces to \(n^2\) and the third weight vanishes, so \(\pi_{r:0\rightarrow r'}=\pi\) as opinion distributions and the numerator vanishes. The numerator is therefore divisible by \(k\), and direct expansion and simplification yield
\[
\Delta(k)=\frac{k\,R_{m,n,d,e}(k)}{Q_{m,n}(k)},
\]
with \(R_{m,n,d,e}(k)\) the cubic polynomial given in the statement.

Finally, completing the square in the inner factor of \(Q_{m,n}\),
\[
2k^2-2nk+C
=2\bigl(k-\tfrac{n}{2}\bigr)^{2}+\bigl(m^2+\tfrac{n^2}{2}\bigr)>0
\quad\text{for all }k\in\mathbb{R}.
\]
Since also \(C=m^2+n^2>0\), we have \(Q_{m,n}(k)>0\) on the whole real line, in particular on \([0,n]\). Therefore, for \(k>0\), the sign of \(\Delta(k)\) is determined by \(R_{m,n,d,e}(k)\).
\end{proof}

\begin{proposition}[Uniqueness of the tipping point]
\label{prop:unique-root-tipping-general}
Let \(m,n,d,e>0\). Then the polynomial \(R_{m,n,d,e}\) has exactly one root in \((0,n)\). Moreover, this root is simple.
\end{proposition}

\begin{proof}
Set \(x\defsymbol n^2/m^2>0\), \(s\defsymbol e/d>0\), and \(t\defsymbol k/n\). Then \(k\in(0,n)\) if and only if \(t\in(0,1)\), and a direct computation gives
\[
R_{m,n,d,e}(k)=d^2m^4n\,\widetilde R_{x,s}(t),
\]
where
\[
\widetilde R_{x,s}(t)
=
s^2x(x+1)t^3-2s^2x(x+1)t^2+\bigl(s^2(x+1)^2+2s(x+1)+2\bigr)t-2.
\]
Since \(d^2m^4n>0\), the roots of \(R_{m,n,d,e}\) in \((0,n)\) are exactly the roots of \(\widetilde R_{x,s}\) in \((0,1)\).

Now \(\widetilde R_{x,s}(0)=-2<0\), while
\[
\widetilde R_{x,s}(1)=s(x+1)(s+2)>0.
\]
Hence \(\widetilde R_{x,s}\) has at least one root in \((0,1)\).

To prove uniqueness, compute the discriminant of \(\widetilde R_{x,s}\) with respect to \(t\):
\[
\Disc_t(\widetilde R_{x,s})=-4s^2x(x+1)\,P(x,s),
\]
where
\[
P(x,s)=(s^6+2s^5)x^5+x^2(Ax^2+Bx+C)+Dx+E,
\]
with
\[
\begin{aligned}
A&=5s^6+14s^5+8s^4,\\
B&=10s^6+36s^5+26s^4-12s^3,\\
C&=10s^6+44s^5+46s^4+8s^3+23s^2,
\end{aligned}
\]
and
\[
\begin{aligned}
D&=5s^6+26s^5+46s^4+52s^3+59s^2+24s,\\
E&=s^6+6s^5+18s^4+32s^3+36s^2+24s+8.
\end{aligned}
\]
Clearly \(A,C,D,E>0\). Moreover,
\[
4AC-B^2
=
4s^6(s+2)\bigl(25s^5+130s^4+212s^3+204s^2+234s+74\bigr)>0.
\]
Thus the quadratic polynomial \(Ax^2+Bx+C\) is positive for every \(x>0\). Since also \((s^6+2s^5)x^5>0\), \(Dx>0\), and \(E>0\), it follows that \(P(x,s)>0\) for all \(x,s>0\). Therefore
\[
\Disc_t(\widetilde R_{x,s})<0.
\]
A cubic with negative discriminant has exactly one real root. Hence \(\widetilde R_{x,s}\) has exactly one real root. Since \(\widetilde R_{x,s}(0)<0<\widetilde R_{x,s}(1)\), that root lies in \((0,1)\). Therefore \(R_{m,n,d,e}\) has exactly one root \(k^\star\in(0,n)\). Since the discriminant is non-zero, this root is simple.
\end{proof}

\subsection*{Proof of Proposition~\ref{prop:mec22-mass-transfer}}

\massTransfer*
\begin{proof}
Let \(w_i=p_i^2\), \(D=\sum_{i=1}^r w_i\), and \(N=\sum_{1\leq i<j\leq r}w_iw_j(x_i-x_j)^2\). By Lem.~\ref{lem:weighted-variance},
\[
\MEC[2,2](\pi)=\frac{N}{D}.
\]
For \(\pi^{(k)}\), set \(x_{r+1}=1\) and
\[
w_i(k)=
\begin{cases}
p_i^2 & i\le r-1,\\
(p_r-k)^2 & i=r,\\
k^2 & i=r+1,
\end{cases}
\]
and let \(D^{(k)}=\sum_{i=1}^{r+1}w_i(k)=\sum_{i=1}^{r-1}p_i^2+(p_r-k)^2+k^2\) and \(N^{(k)}=\sum_{1\leq i<j\leq r+1}w_i(k)w_j(k)(x_i-x_j)^2\). By Lem.~\ref{lem:weighted-variance} again, \(\MEC[2,2](\pi^{(k)})=N^{(k)}/D^{(k)}\).

Set \(Q_\pi(k)\defsymbol D\,D^{(k)}\). Since \(a^2+b^2\geq\tfrac12(a+b)^2\),
\[
D^{(k)}\geq(p_r-k)^2+k^2\geq\tfrac12 p_r^2>0,
\]
so \(Q_\pi(k)>0\) for every \(k\in\mathbb{R}\), in particular on \([0,p_r]\).

Putting \(\MEC[2,2](\pi^{(k)})-\MEC[2,2](\pi)\) over the common denominator \(Q_\pi(k)\),
\[
\MEC[2,2](\pi^{(k)})-\MEC[2,2](\pi)=\frac{N^{(k)}D-N\,D^{(k)}}{Q_\pi(k)}.
\]
The weights \(w_i(k)\) are polynomials in \(k\) of degree at most \(2\), so \(N^{(k)}\) is a polynomial of degree at most \(4\) and \(D^{(k)}\) a polynomial of degree \(2\). Hence the numerator above is a polynomial of degree at most \(4\). At \(k=0\) we have \(w_{r+1}(0)=0\) and \(w_r(0)=p_r^2\), so \(D^{(0)}=D\) and \(N^{(0)}=N\), giving \(N^{(0)}D-N\,D^{(0)}=0\). The numerator is therefore divisible by \(k\), and we may write
\[
\MEC[2,2](\pi^{(k)})-\MEC[2,2](\pi)=\frac{k\,P_\pi(k)}{Q_\pi(k)}
\]
with \(P_\pi(k)\) a polynomial of degree at most \(3\).

Finally, \(P_\pi\) is genuinely cubic. The only contribution to \(k^4\) in \(N^{(k)}D-N\,D^{(k)}\) comes from the cross term \(w_r(k)w_{r+1}(k)(x_r-1)^2D=(p_r-k)^2k^2(1-x_r)^2D\), whose \(k^4\) coefficient is \(D(1-x_r)^2\). Since \(x_r<1\) and \(D>0\),
\[
[k^3]P_\pi(k)=D\,(1-x_r)^2>0,
\]
so \(P_\pi\) has degree exactly \(3\) with positive leading coefficient.
\end{proof}